\section*{الفصل \ref{ch:g5:decimals} --- الأعداد العشرية}

\begin{solution}{exo:g5:decimals:1}
$5.207$; \quad $0.85$; \quad $4.507$; \quad $12.009$.
\end{solution}

\begin{solution}{exo:g5:decimals:2}
الرقم $4$ يعدّ الأعشار، والرقم $8$ يعدّ الأجزاء من ألف:
\[
27.418 = 27 + \frac{4}{10} + \frac{1}{100} + \frac{8}{1000}.
\]
\end{solution}

\begin{solution}{exo:g5:decimals:3}
$3.60 = 3.6$: صواب (\omterm{def:g1:counting:zero}{فالصفر} النهائي لا يضيف شيئًا).

$0.5 = 0.05$: خطأ --- $5$ أعشار مقابل $5$ أجزاء من مئة.

$2.070 = 2.07$: صواب.

$1.3 = 1.03$: خطأ --- $3$ أعشار مقابل $3$ أجزاء من مئة.
\end{solution}

\begin{solution}{exo:g5:decimals:4}
$4.8 > 4.79$؛ \quad $0.302 < 0.32$ (الأعشار متساوية، والأجزاء من مئة
$0 < 2$)؛ \quad $6.5 = 6.500$؛ \quad $9.09 < 9.1$.
\end{solution}

\begin{solution}{exo:g5:decimals:5}
$3.041 < 3.104 < 3.14 < 3.4 < 3.41$.
\end{solution}

\begin{solution}{exo:g5:decimals:6}
$6.83$: بين $6$ و $7$؛ وبين $6.8$ و $6.9$.

$0.492$: بين $0$ و $1$؛ وبين $0.4$ و $0.5$.

$12.06$: بين $12$ و $13$؛ وبين $12.0$ و $12.1$.
\end{solution}

\begin{solution}{exo:g5:decimals:7}
$8.276$: إلى الوحدة $8$؛ وإلى العُشر $8.3$؛ وإلى الجزء من مئة
$8.28$. وأما $3.95$ إلى العُشر: فرقم الأجزاء من مئة هو $5$، فنقرّب
إلى الأعلى: $4.0$.
\end{solution}

\begin{solution}{exo:g5:decimals:8}
مثلًا $2.65$؛ \quad $0.415$؛ \quad $5.995$. (وهناك أجوبة لا
تُحصى في كل مرة.)
\end{solution}

\begin{solution}{exo:g5:decimals:9}
$0.85 < 0.9 < 1.02 < 1.2$ (بالأمتار).
\end{solution}

\begin{solution}{exo:g5:decimals:10}
قارن الأعشار أولًا: في $7.12$ عُشر واحد ($1$)، وفي $7.9$ تسعة
أعشار ($9$)، إذن $7.12 < 7.9$. وقد قارنت لُبنى الجزأين العشريين
كعددين صحيحين ($12$ مقابل $9$)، ناسيةً أن $12$ جزءًا من مئة أقلّ
بكثير من $90$ جزءًا من مئة.
\end{solution}

\begin{solution}{exo:g5:decimals:11}
\omterm{def:g4:numbers:round}{التقريب} إلى العُشر يعطي $6.4$ للأعداد ذات الرقمين العشريين من
$6.35$ (فالمنتصف يقرَّب إلى الأعلى) إلى $6.44$: أي الأعداد $6.35$ و $6.36$
و \dots و $6.44$. فالأصغر: $6.35$؛ والأكبر: $6.44$.
\end{solution}
